\documentclass[../../../include/open-logic-section]{subfiles}

\begin{document}

\olfileid{sth}{spine}{recursion}
\olsection{અનંતાતીત પુનરાવર્તનનાં પ્રમેયો}

જોકે સૌથી પહેલાં આપણે ખાતરી કરવી પડશે કે
\olref[valpha]{defValphas} સફળ વ્યાખ્યા છે. વધુ સામાન્ય રીતે,
અનંતાતીત પુનરાવર્તનથી વ્યાખ્યા આપવાનો દરેક પ્રયત્ન સફળ થશે તે
સાબિત કરવું પડશે. આ વિભાગનો એ જ હેતુ છે.

\emph{સાવધાન: અહીંનો વિષય અટપટો છે}. છતાં મુખ્ય સાર બહુ સરળ છે:
અનંતાતીત અનુમાન અને પ્રતિસ્થાપન મળીને અનંતાતીત પુનરાવર્તનનાં
ઘણાં સ્વરૂપોને વાજબી ઠેરવે છે.\footnote{યાદ રાખો: જુદું સ્પષ્ટ ન
કર્યું હોય તો બધાં સૂત્રો અને પદોમાં પરિમાણો હોઈ શકે છે.}

\begin{defn}
	ધારો કે $\tau(x)$ પદ, $f$ વિધેય અને $\alpha$ ક્રમવાચક સંખ્યા છે.
	જો $\dom{f} = \alpha$ અને $(\forall \beta \in \alpha)f(\beta) = \tau(\funrestrictionto{f}{\beta})$ બંને હોય,
	તો અને માત્ર તો $f$ને $\tau$ માટેનું \emph{$\alpha$-સન્નિકટન} કહીએ.
\end{defn}

\begin{lem}[મર્યાદિત પુનરાવર્તન]\ollabel{transrecursionfun}
કોઈ પણ પદ $\tau(x)$ અને ક્રમવાચક સંખ્યા $\alpha$ માટે,
$\tau$નું એક અનન્ય $\alpha$-સન્નિકટન છે.
\end{lem}
\begin{proof}
અમે બતાવીશું કે કોઈ પણ $\gamma \leq \alpha$ માટે અનન્ય
$\gamma$-સન્નિકટન છે.

પહેલાં અનન્યતા સ્થાપીએ. ધારો કે $g$ અને $h$ અનુક્રમે
$\gamma$- અને $\delta$-સન્નિકટનો છે. તેમની દલીલો પર
અનંતાતીત અનુમાનથી બતાવી શકાય કે કોઈ પણ
$\beta \in \dom{g} \cap \dom{h} =
\gamma \cap \delta = \min(\gamma, \delta)$ માટે
$g(\beta) = h(\beta)$. તેથી સન્નિકટનો અસ્તિત્વમાં હોય તો
અનન્ય છે અને સમાન દલીલો પર તેમનાં મૂલ્યો સરખાં છે.

અસ્તિત્વ સ્થાપવા માટે હવે $\delta \leq \alpha$ ક્રમવાચક
સંખ્યાઓ પર સરળ અનંતાતીત અનુમાન
(\olref[ordinals][opps]{simpletransrecursion}) વાપરીએ.

ખાલી વિધેય સ્પષ્ટપણે $\emptyset$-સન્નિકટન છે.

જો $g$ એ $\gamma$-સન્નિકટન હોય, તો $g \cup
\{\tuple{\gamma, \tau(g)}\}$ એ $\ordsucc{\gamma}$-સન્નિકટન છે.

જો $\gamma$ સીમા ક્રમવાચક સંખ્યા હોય અને દરેક $\delta < \gamma$
માટે $g_\delta$ એ $\delta$-સન્નિકટન હોય, તો
$g = \bigcup_{\delta \in \gamma} g_\delta$ લો. જુદાં જુદાં
$g_\delta$ના મૂલ્યો જ્યાં મળે ત્યાં સરખાં હોવાથી આ વિધેય છે.
વળી $\delta \in \gamma$ હોય, તો
$g(\delta) = g_{\ordsucc{\delta}}(\delta) =
\tau(\funrestrictionto{g_{\ordsucc{\delta}}}{\delta}) =
\tau(\funrestrictionto{g}{\delta})$.

આથી અનંતાતીત અનુમાન વડે સાબિતી પૂરી થાય છે.
\end{proof}

જો વિધેયને બદલે \emph{પદ} વ્યાખ્યાયિત કરવાની છૂટ લઈએ, તો અગાઉના
પરિણામમાંથી મર્યાદા $\alpha$ દૂર કરી શકીએ. આગળના પરિણામના નિવેદન
અને પુરાવામાં $\sigma$ પદ હોય ત્યારે
$\funrestrictionto{\sigma}{\alpha} = \Setabs{\tuple{\beta,
\sigma(\beta)}}{\beta \in \alpha}$ લઈશું.

\begin{thm}[સામાન્ય પુનરાવર્તન]\ollabel{transrecursionschema}
કોઈ પણ પદ $\tau(x)$ માટે આપણે સ્પષ્ટ રીતે એવું પદ
$\sigma(x)$ વ્યાખ્યાયિત કરી શકીએ કે કોઈ પણ ક્રમવાચક
સંખ્યા~$\alpha$ માટે
$\sigma(\alpha) = \tau(\funrestrictionto{\sigma}{\alpha})$.
\end{thm}

\begin{proof}
દરેક $\alpha$ માટે \olref{transrecursionfun} મુજબ $\tau$નું
અનન્ય $\alpha$-સન્નિકટન $f_\alpha$ છે.
$\sigma(\alpha)$ને $f_{\ordsucc{\alpha}}(\alpha)$ તરીકે વ્યાખ્યાયિત
કરો. હવે:
	\begin{align*}
		\sigma(\alpha) &= 
		f_{\ordsucc{\alpha}}(\alpha) \\&= 
		\tau(\funrestrictionto{f_{\ordsucc{\alpha}}}{\alpha}) \\&= 
		\tau(\Setabs{\tuple{\beta, f_{\ordsucc{\alpha}}(\beta)}}{\beta \in \alpha}) \\&= 
		\tau(\Setabs{\tuple{\beta, f_{\ordsucc{\beta}}(\beta)}}{\beta \in \alpha})\\&=
		\tau(\funrestrictionto{\sigma}{\alpha})
	\end{align*}
અહીં \olref{transrecursionfun} મુજબ બધા $\beta < \alpha$ માટે
$f_{\ordsucc{\beta}}(\beta) = f_{\ordsucc{\alpha}}(\beta)$ છે.
\end{proof}
\noindent 
ધ્યાન આપો કે \olref{transrecursionschema} એક \emph{પ્રરૂપ} છે.
ખાસ કરીને, $\sigma$ વિધેય, એટલે કે ચોક્કસ પ્રકારનો \emph{ગણ},
વ્યાખ્યાયિત કરે એવી અપેક્ષા ન રાખી શકાય. એમ હોય તો
$\dom{\sigma}$ બધી ક્રમવાચક સંખ્યાઓનો ગણ થાય, જે
બુરાલી–ફૉર્ટીના વિરોધાભાસ
(\olref[ordinals][basic]{buraliforti})ની વિરુદ્ધ છે.

છતાં \olref{transrecursionschema}થી $V_\alpha$ની આપણી વ્યાખ્યા
વાજબી ઠરે છે તે બતાવવાનું બાકી છે. તે તરત સ્પષ્ટ નહીં થાય;
પરંતુ અનંતાતીત પુનરાવર્તનના એક અંતિમ સરળ સ્વરૂપથી સમજાશે.

\begin{thm}[સરળ પુનરાવર્તન]\ollabel{simplerecursionschema} 
કોઈ પણ પદો $\tau(x)$ અને $\theta(x)$ તથા કોઈ પણ ગણ $A$ માટે
આપણે સ્પષ્ટ રીતે એવું પદ $\sigma(x)$ વ્યાખ્યાયિત કરી શકીએ કે:
\begin{align*}
	\sigma(\emptyset) &= A\\
	\sigma(\ordsucc{\alpha}) &= \tau(\sigma(\alpha)) &&
		\text{for any ordinal }\alpha\\
	\sigma(\alpha) &= \theta(\ran{\funrestrictionto{\sigma}{\alpha}})&&
	\text{when }\alpha\text{ is a limit ordinal}
\end{align*}
\end{thm}

\begin{proof}
પહેલાં $\xi(x)$ પદને આ રીતે વ્યાખ્યાયિત કરીએ:
%t $\xi(x) = A$ if $x$ is not a function whose domain is an ordinal;
%otherwise let:
\[
	\xi(x) = 
	\begin{cases}
			A &\text{if $x$ is not a function whose}\\
			&\hspace{1em}\text{domain is an ordinal; otherwise:}\\
			\tau(x(\alpha)) & \text{if $\dom{x} = \ordsucc{\alpha}$}\\
			\theta(\ran{x}) & \text{if $\dom{x}$ is a limit ordinal}
	\end{cases}
\]
\sourcecorrection{OLMD-118}{મૂળની ખંડવાર વ્યાખ્યામાં ખાલી વિધેય
છૂટી જાય છે: તે વિધેય છે અને તેનો પ્રદેશ ખાલી ક્રમવાચક સંખ્યા છે,
એટલે તે પહેલી શરત સંતોષતો નથી; તેનો પ્રદેશ અનુગામી કે સીમા
ક્રમવાચક પણ નથી. પહેલી શરતમાં ખાલી વિધેયને પણ સમાવીને
તેના માટે $A$ મૂલ્ય લેવું જરૂરી છે. તેથી જ આગળ
$\xi(\emptyset) = A$ વાપરી શકાય છે.}
\olref{transrecursionschema} મુજબ એવું પદ $\sigma(x)$ છે કે
દરેક ક્રમવાચક સંખ્યા $\alpha$ માટે
$\sigma(\alpha) = \xi(\funrestrictionto{\sigma}{\alpha})$;
વળી $\funrestrictionto{\sigma}{\alpha}$ એ પ્રદેશ $\alpha$
ધરાવતું વિધેય છે. સરળ અનંતાતીત અનુમાનથી
(\olref[ordinals][opps]{simpletransrecursion}) બતાવીશું કે
$\sigma$માં જરૂરી ગુણધર્મો છે.

પહેલાં, $\sigma(\emptyset) = \xi(\emptyset) = A$.

પછી, $\sigma(\ordsucc{\alpha}) = \xi(\funrestrictionto{\sigma}{\ordsucc{\alpha}}) = \tau(\funrestrictionto{\sigma}{\ordsucc{\alpha}}(\alpha)) = \tau(\sigma(\alpha))$.

છેલ્લે, $\alpha$ સીમા ક્રમવાચક સંખ્યા હોય ત્યારે
$\sigma(\alpha) = \xi(\funrestrictionto{\sigma}{\alpha}) = \theta(\ran{\funrestrictionto{\sigma}{\alpha}})$.
\end{proof}
\noindent
હવે \olref[valpha]{defValphas}ને વાજબી ઠેરવવા માત્ર $A
= \emptyset$, $\tau(x) = \Pow{x}$ અને $\theta(x) = \bigcup x$
લો. આથી છેવટે $V_\alpha$ની વ્યાખ્યા વાજબી ઠરે છે!{}


\end{document}
